% =====================================================================
%  2bat-ud2-1.tex  —  SESSIÓ 1: Les inequacions de 1r, ara amb dues variables
%  (sense preàmbul: s'inclou des de main.tex)
% =====================================================================
\horatitol{Sessió 1 — Les inequacions de 1r, ara amb dues variables}%
{Recuperar les inequacions d'una variable i fer el salt al pla: quan hi ha dues quantitats que varien alhora.}

\subsection{Idea clau}

A 1r vam treballar inequacions d'\textbf{una} variable (llindars de renda, franges
d'edat) i les representàvem sobre la \textbf{recta}. Aquest curs sovint haurem de
decidir sobre \textbf{dues} quantitats alhora: per exemple, com repartir uns diners
entre \emph{dos} programes. Quan hi ha dues variables, una recta ja no n'hi ha prou:
necessitem tot el \textbf{pla}.

\begin{idea}
\textbf{D'una variable a dues}
\begin{itemize}[leftmargin=1.4em, itemsep=2pt]
  \item Amb \textbf{una} variable, una inequació com $x\le 5$ té per solució un
        \textbf{tram de la recta}.
  \item Amb \textbf{dues} variables ($x$ i $y$), una inequació com $x+y\le 10$ té per
        solució tota una \textbf{zona del pla}: totes les parelles $(x,y)$ que la
        compleixen.
  \item Representem les parelles $(x,y)$ com a \textbf{punts} en uns eixos: $x$
        horitzontal, $y$ vertical.
\end{itemize}
\textbf{Exemple social:} si una entitat reparteix diners entre el programa A ($x$) i
el programa B ($y$) i el total no pot superar $\num{10000}\eur$, la condició és
$x+y\le\num{10000}$. Hi ha \emph{moltes} maneres de complir-la.
\end{idea}

\subsection{Exemples}

\begin{exemple}[repàs: una inequació d'una variable]
Recordem 1r. «Una entitat té un pressupost que no pot superar $\num{8000}\eur$»: si
$x$ és la despesa, $x\le\num{8000}$, l'interval $(-\infty,\,\num{8000}]$. A la recta,
un punt ple a $\num{8000}$ i tot el que queda a l'esquerra. \textbf{Una} sola
quantitat, \textbf{una} recta.
\end{exemple}

\begin{exemple}[el salt al pla: dues variables]
Ara l'entitat reparteix els diners entre \emph{dos} programes: $x\eur$ al programa A
i $y\eur$ al programa B, amb un total que no pot superar $\num{10000}\eur$:
\[
x+y\le \num{10000}.
\]
Algunes combinacions que la compleixen: $(\num{5000},\num{4000})$ ja que
$\num{5000}+\num{4000}=\num{9000}\le\num{10000}$; o $(\num{2000},\num{1000})$. Una que
\emph{no} la compleix: $(\num{7000},\num{6000})$, perquè $\num{13000}>\num{10000}$.
Cada combinació és un \textbf{punt} $(x,y)$; les que compleixen la condició formen una
\textbf{zona} del pla:

\begin{center}
\begin{tikzpicture}
\begin{axis}[udaxis, width=8.6cm, height=6cm,
    xlabel={\small programa A ($x$, milers \euro)},
    ylabel={\small programa B ($y$, milers \euro)},
    xmin=0, xmax=12, ymin=0, ymax=12,
    xtick={0,5,10}, ytick={0,5,10}, clip=false]
  \addplot[name path=recta, udblue, domain=0:10, samples=2]{10-x};
  \path[name path=eixx] (0,0) -- (10,0);
  \addplot[udblue!18] fill between[of=recta and eixx, soft clip={domain=0:10}];
  \addplot[udnavy, only marks, mark=*, mark size=1.6pt] coordinates {(5,4)};
  \node[udnavy, font=\scriptsize, anchor=south west] at (axis cs:5,4) {$(5,4)$};
  \node[udblue, font=\scriptsize, anchor=west] at (axis cs:1.2,2) {compleixen};
\end{axis}
\end{tikzpicture}

\figcap{Cada punt és un repartiment $(x,y)$; la zona ombrejada són els que compleixen $x+y\le 10$ (en milers).}
\end{center}
\end{exemple}

\subsection{Exercicis resolts}

\begin{exresolt}[comprovar si un punt compleix una condició]
Per a la condició $x+y\le\num{10000}$ (en euros), digues si compleixen cada
repartiment: (a) $(\num{6000},\num{3000})$; (b) $(\num{8000},\num{5000})$;
(c) $(\num{4000},\num{6000})$.
\solucio
Substituïm i sumem:
\begin{enumerate}[label=\alph*), leftmargin=1.6em]
  \item $\num{6000}+\num{3000}=\num{9000}\le\num{10000}$: \textbf{sí}.
  \item $\num{8000}+\num{5000}=\num{13000}\le\num{10000}$ és fals: \textbf{no}.
  \item $\num{4000}+\num{6000}=\num{10000}\le\num{10000}$: \textbf{sí} (just al límit).
\end{enumerate}
\resposta{(a) sí; (b) no; (c) sí (just al límit, que hi entra).}
\end{exresolt}

\begin{exresolt}[traduir una frase a una inequació de dues variables]
Una associació dedica $x$ hores setmanals a tallers i $y$ hores a acompanyament, i en
total disposa com a màxim de $30$ hores. A més, totes dues han de ser no negatives.
Escriu les condicions com a inequacions.
\solucio
«En total com a màxim $30$ hores»: $x+y\le 30$. «No negatives» (no es poden fer hores
negatives): $x\ge 0$ i $y\ge 0$. Així, el conjunt de condicions és
\[
x+y\le 30,\qquad x\ge 0,\qquad y\ge 0.
\]
\resposta{$x+y\le 30$, amb $x\ge 0$ i $y\ge 0$.}
\end{exresolt}

\subsection{Exercicis per fer}

\begin{exercicis}
\begin{exlist}
  \item Resol aquestes inequacions d'\textbf{una} variable (repàs de 1r) i escriu-ne
        l'interval:
        \begin{enumerate}[label=\alph*), leftmargin=1.6em, topsep=2pt]
          \item $2x-4\le 10$ \hfill \item $5-3x>-1$
        \end{enumerate}
  \item Per a la condició $x+y\le 12$, digues si compleixen els punts:
        \begin{enumerate}[label=\alph*), leftmargin=1.6em, topsep=2pt]
          \item $(5,4)$ \hfill \item $(8,6)$ \hfill \item $(7,5)$
        \end{enumerate}
  \item Tradueix cada frase a una inequació amb dues variables ($x$ i $y$):
        \begin{enumerate}[label=\alph*), leftmargin=1.6em, topsep=2pt]
          \item «El cost del programa A més el del B no pot superar $\num{5000}\eur$.»
          \item «Entre tallers ($x$) i sortides ($y$) es faran com a mínim $8$ activitats.»
        \end{enumerate}
  \item Una entitat reparteix $x$ places al servei de matí i $y$ a la tarda, amb un
        total màxim de $50$ places i totes dues no negatives. Escriu les tres
        condicions (la del total i les dues de no-negativitat).
  \item Dóna \textbf{tres} repartiments diferents $(x,y)$ que compleixin
        $x+y\le 20$ amb $x\ge 0$ i $y\ge 0$.
  \item \textbf{(Reflexió)} Explica amb les teves paraules què canvia quan, en lloc
        d'\emph{una} incògnita, en tenim \emph{dues}. Per què ja no n'hi ha prou amb
        una recta per representar la solució?
\end{exlist}
\end{exercicis}

% ---------------------------------------------------------------------
\fullsolucions{Sessió 1}
\begin{solucions}
\begin{sollist}
  \item (a) $2x\le 14 \Rightarrow x\le 7$, interval $(-\infty,7]$; \quad
        (b) $-3x>-6 \Rightarrow x<2$ (girant el signe), interval $(-\infty,2)$.
  \item (a) $5+4=9\le 12$: sí; \quad (b) $8+6=14\le 12$ fals: no; \quad
        (c) $7+5=12\le 12$: sí (just al límit).
  \item (a) $x+y\le\num{5000}$; \quad (b) $x+y\ge 8$.
  \item $x+y\le 50$, amb $x\ge 0$ i $y\ge 0$.
  \item Resposta oberta. Exemples vàlids: $(10,5)$, $(0,20)$, $(8,8)$ (tots compleixen
        $x+y\le 20$ amb $x,y\ge 0$).
  \item Resposta oberta. Idea esperada: amb dues variables, una solució és una
        \emph{parella} $(x,y)$, no un sol nombre; el conjunt de solucions ja no és un
        tram de recta sinó una \textbf{zona del pla}, i per això cal representar-ho en
        dos eixos.
\end{sollist}
\end{solucions}
